\scanpage{75}
\uline{Remarque} Il \uline{n'est pas vrai} que la seule donnée des $\widetilde A_{M,i}$ permet de retrouver \deleted{les} $\widetilde F_{M,i}$, par la \unsure{récette|recette} naïve --- car \inserted{pour} un $f\in\mathfrak P_3(\widetilde S_M)$, $f=\{x,y,z\}$, il ne suffit pas qu'on ait $a\in\widetilde A_{M,i}\ \forall a\in\mathfrak P_2(f)$, i.e $x\wedge y,y\wedge z,z\wedge x\in i=\{\pm e\}\subset\det M^*$, pour avoir $f\in\widetilde F_{M,i}$ i.e. $x+y+z=0$. En effet, on constate que la condition envisagée signifie simplement que $z=\pm x\pm y$ (et non $z=-(x+y)$.
Mais dans le cas $2.1_\Lambda=0$ --- qui est celui où c'est la structure $\widetilde P_M$ \deleted{\unsure{qui}}\inserted{elle-même} \unsure{est|où} la plus intéressante --- cela implique bien $x+y+z=0$, i.e.
\[
 \text{Si }2.1_\Lambda=0,\quad\text{alors }\widetilde F_{M,i}=\bigl\{f\in\mathfrak P_3(\widetilde S_M)\mid
 \substack{\forall a\in\mathfrak P_2(f),\\\text{on a }a\in\widetilde A_{M,i}}\bigr\}.\tag*{(15)}
\]

\uline{Question de connexité des} $\widetilde P_{M,i}$. On s'y intéresse, car on aimerait qu'un autom. de $\widetilde P_{M,i}$ (i.e. permutation de $\widetilde S_M$ qui \deleted{\unsure{applique}}\inserted{transforme} $\widetilde F_{M,i}$ en lui-\unsure{$=$|même}) qui fixe un \uline{repère} $(x,y,z)$ (i.e. un $(x,y,z)\in\widetilde S_M^3$ tels que $\{x,y,z\}\in\widetilde F_{M,i}$) soit l'identité.
% slanted
\margin{donc il faut considérer l'antipodisme comme faisant partie de la structure \unsure{enrichie} de $\widetilde P_{M,i}$ \ldots\ '' )!}
[Pour pouvoir travailler, il \unsure{le faut|faut} au moins que l'autom. commute à ``l'antipodisme'' \deleted{\unsure{$\varepsilon$|$\in$}} $\underline{\widetilde a}_M$ (autom de $\widetilde P_{M,i}$ induit par $-\mathrm{id}_M$). Alors, si $u$ fixe un $\unsure{s|\xi|x}\in\widetilde S_M$,
% archive sheet number: 74
